\subsection{Fit time nelle figure del paper}

\reviewcomment{Reviewer 3, comment 3}{
Because this is a submission to Statistics and Computing, I think the computational side should
be developed more explicitly. The paper explains the algorithmic structure and mentions
implementation in an open-source C++ library, which is a real strength. However, the manuscript
currently provides very little quantitative information on computational cost, scalability, or
numerical behavior. This is especially important in the neuroimaging application, where the
predictors live on a mesh with more than 36,000 nodes. At a minimum, I would like to see
\reviewhighlight{additional reporting on runtime and, if possible, memory usage for the main
experiments; number of iterations required for convergence; practical computational cost of
GCV-based tuning; sensitivity of computation to the basis dimension or mesh resolution; and
discussion of numerical stability when extracting several components.} Even a modest runtime
or scaling study would substantially improve the paper for this journal.}

\notebox{Come deciso, nelle figure quantitative già presenti nel paper la legend viene
spostata sotto i pannelli e lo spazio liberato in basso a destra viene usato per
il fit time. La composizione mostrata qui è ancora provvisoria: sistemo le figure
definitive una volta concordati i risultati da inserire. Le tabelle riportano la mediana delle
iterazioni per componente su tutti i livelli e le ripetizioni del test; ``100\%'' indica che tutte
le traiettorie osservate sono monotone. Il trattino indica una diagnostica non disponibile oppure
un fit non iterativo. Il fit time include preprocessing e tuning specifici, fit e post-processing
in memoria; esclude generazione dei dati, serializzazione e caricamento dei risultati, e
valutazione.}

\subsubsection{Fig.~3: \texttt{1D/vary\_k}}
\begin{figure}[H]
\centering
\begin{subfigure}[t]{0.48\linewidth}
\centering\includegraphics[width=\linewidth]{../images/1D/vary_k/Y_reconstruction.pdf}
\caption{RMSE $Y$.}
\end{subfigure}\hfill
\begin{subfigure}[t]{0.48\linewidth}
\centering\includegraphics[width=\linewidth]{../images/1D/vary_k/X_reconstruction_locations.pdf}
\caption{RMSE $X$.}
\end{subfigure}
\par\vspace{0.35em}
\begin{subfigure}[t]{0.48\linewidth}
\centering\includegraphics[width=\linewidth]{../images/1D/vary_k/Beta1_locations.pdf}
\caption{RMSE $\beta_1$.}
\end{subfigure}\hfill
\begin{subfigure}[t]{0.48\linewidth}
\centering\includegraphics[width=\linewidth]{../images/1D/vary_k/fit_seconds.pdf}
\caption{Fit time.}
\end{subfigure}
\par\vspace{0.15em}
\modellegend
\caption{Proposta per la Fig.~3: legend sotto i pannelli, fit time in basso a destra.}
\end{figure}

\reviewiterationtable{review_tables/1D_vary_k_iterations.csv}

\notebox{La figura riporta il fit time al variare di $k$. B-fPLS con smoothing splines è il
metodo più costoso: effettua il pre-smoothing separatamente per ciascuna osservazione e usa una
griglia di $\lambda$ più densa. R1-fPLS e fPCR sono più veloci di PB-fPLS e di B-fPLS con
smoothing splines, con un vantaggio che cresce all'aumentare di $k$; B-fPLS con regression
splines e MV-PLS restano più leggeri in una parte dell'intervallo.}

\subsubsection{Fig.~5: \texttt{1D/vary\_n}}
\begin{figure}[H]
\centering
\begin{subfigure}[t]{0.48\linewidth}
\centering\includegraphics[width=\linewidth]{../images/1D/vary_n/Y_reconstruction.pdf}
\caption{RMSE $Y$.}
\end{subfigure}\hfill
\begin{subfigure}[t]{0.48\linewidth}
\centering\includegraphics[width=\linewidth]{../images/1D/vary_n/X_reconstruction_locations.pdf}
\caption{RMSE $X$.}
\end{subfigure}
\par\vspace{0.35em}
\begin{subfigure}[t]{0.48\linewidth}
\centering\includegraphics[width=\linewidth]{../images/1D/vary_n/Beta1_locations.pdf}
\caption{RMSE $\beta_1$.}
\end{subfigure}\hfill
\begin{subfigure}[t]{0.48\linewidth}
\centering\includegraphics[width=\linewidth]{../images/1D/vary_n/fit_seconds.pdf}
\caption{Fit time.}
\end{subfigure}
\par\vspace{0.15em}
\modellegend
\caption{Proposta per la Fig.~5: legend sotto i pannelli, fit time in basso a destra.}
\end{figure}

\reviewiterationtable{review_tables/1D_vary_n_iterations.csv}

\notebox{Qui R1-fPLS e fPCR usano la discretizzazione con 51 nodi, mentre B-fPLS e
PB-fPLS usano 11 nodi. R1-fPLS rimane comunque più veloce di PB-fPLS e di B-fPLS con
smoothing splines, mentre B-fPLS con regression splines e MV-PLS restano più leggeri in
questo setting.}

\subsubsection{Fig.~S2: \texttt{1D/vary\_m}}
\begin{figure}[H]
\centering
\begin{subfigure}[t]{0.48\linewidth}
\centering\includegraphics[width=\linewidth]{../images/1D/vary_m/Y_reconstruction.pdf}
\caption{RMSE $Y$.}
\end{subfigure}\hfill
\begin{subfigure}[t]{0.48\linewidth}
\centering\includegraphics[width=\linewidth]{../images/1D/vary_m/X_reconstruction_locations.pdf}
\caption{RMSE $X$.}
\end{subfigure}
\par\vspace{0.35em}
\begin{subfigure}[t]{0.48\linewidth}
\centering\includegraphics[width=\linewidth]{../images/1D/vary_m/Beta1_locations.pdf}
\caption{RMSE $\beta_1$.}
\end{subfigure}\hfill
\begin{subfigure}[t]{0.48\linewidth}
\centering\includegraphics[width=\linewidth]{../images/1D/vary_m/fit_seconds.pdf}
\caption{Fit time.}
\end{subfigure}
\par\vspace{0.15em}
\modellegend
\caption{Proposta per la Fig.~S2: legend sotto i pannelli, fit time in basso a destra.}
\end{figure}

\reviewiterationtable{review_tables/1D_vary_m_iterations.csv}

\notebox{Anche in questo setting R1-fPLS e fPCR usano la discretizzazione con 51 nodi,
mentre B-fPLS e PB-fPLS usano 11 nodi. R1-fPLS è più veloce di PB-fPLS e di B-fPLS con
smoothing splines, mentre B-fPLS con regression splines e MV-PLS restano più leggeri.}

\subsubsection{Fig.~7: \texttt{2D/vary\_n}}
\begin{figure}[H]
\centering
\begin{subfigure}[t]{0.48\linewidth}
\centering\includegraphics[width=\linewidth]{../images/2D/vary_n/X_reconstruction_locations.pdf}
\caption{RMSE $X$.}
\end{subfigure}\hfill
\begin{subfigure}[t]{0.48\linewidth}
\centering\includegraphics[width=\linewidth]{../images/2D/vary_n/Beta1_locations.pdf}
\caption{RMSE $\beta_1$.}
\end{subfigure}
\par\vspace{0.35em}
\begin{subfigure}[t]{0.48\linewidth}
\centering\includegraphics[width=\linewidth]{../images/2D/vary_n/Beta2_locations.pdf}
\caption{RMSE $\beta_2$.}
\end{subfigure}\hfill
\begin{subfigure}[t]{0.48\linewidth}
\centering\includegraphics[width=\linewidth]{../images/2D/vary_n/fit_seconds.pdf}
\caption{Fit time.}
\end{subfigure}
\par\vspace{0.15em}
\modellegendtwod
\caption{Proposta per la Fig.~7: legend sotto i pannelli, fit time in basso a destra.}
\end{figure}

\reviewiterationtable{review_tables/2D_vary_n_iterations.csv}

\notebox{Nel confronto 2D il fit time di fPCR è inferiore a quello di R1-fPLS nel perimetro
misurato. fPCR stima i loading, mentre R1-fPLS stima anche le direzioni e seleziona via GCV un
$\lambda$ per entrambe le quantità a ogni componente.}
